\scanpage{71}
% archive foliation: 70
$X$ \unsure{projective|projectif} lisse \inserted{géo. connexe} sur $k$ corps fini à $q$ éléments.
Par passage à la limite sur les \deleted{\unsure{quotients}} \inserted{coefficients} finis, on trouve
\diagram{\begin{tikzcd}[column sep=small]
0\arrow[r]&H^{i-1}(\overline{X},\mathbf{Q}_\ell(j))_\pi\arrow[r]&H^i(X\deleted{^{\unsure{\circ}}},\mathbf{Q}_\ell(j))\arrow[r]&H^i(\overline{X},\mathbf{Q}_\ell(j))^\pi_{\deleted{\unsure{\pi}}}\arrow[r]&0
\end{tikzcd}}
d'où le résultat
\[
\text{si }i\ne2j,2j+1,\qquad H^i(X,\mathbf{Q}_\ell(j))=0
\]
si $i=2j$
\diagram{\begin{tikzcd}
H^{2j}(X,\mathbf{Q}_\ell(j))={}&H^{2j}(\overline{X},\mathbf{Q}_\ell(j))^\pi\arrow[d,phantom,"\wr" description,"\text{Conj. de Tate}"]\\
&A^j(X/k)\otimes_{\unsure{\mathbf{Z}|\mathbf{Q}}}\mathbf{Q}_\ell
\end{tikzcd}}
si $i=2j+1$
\diagram{\begin{tikzcd}
H^{2j+1}(X,\mathbf{Q}_\ell(j))={}&H^{2j\deleted{\unsure{+1}}}(\overline{X},\mathbf{Q}_\ell(j))_\pi\arrow[d,phantom,"\wr" description,"\substack{\text{Tate + action}\\\text{semi-simple}}"]\\
&A^j(X/k)\otimes_{\unsure{\mathbf{Z}|\mathbf{Q}}}\mathbf{Q}_\ell
\end{tikzcd}}

Soit maintenant $k_n$ l'extension de degré $n$ de $k$ dans $\overline{k}$, et étudions
\[
\varinjlim_n H^i(X_{k_n},\mathbf{Q}_\ell(j))
\]
on trouve $0$ pour $i\ne2j,2j+1$ ; pour $i=2j$, on trouve
\[
H^{2j}(\overline{X},\mathbf{Q}_\ell(j))^{\unsure{\mathfrak{g}|0}}=A^j(\overline{X})\otimes_{\unsure{\mathbf{Z}|\mathbf{Q}}}\mathbf{Q}_\ell,
\]
enfin pour $i=2j+1$, on trouve également $A^j(\overline{X})\otimes_{\mathbf{Z}}\mathbf{Q}_\ell$.
Donc on \unsure{trouve} \inserted{\unsure{du pt de vue motifs}}
\[
\begin{array}{rcl}
\varinjlim H^i(X_{k_n},\mathbf{Q}(j))&=&\begin{cases}
0&\text{si }i\ne2j,2j+1\\
A^j(\overline{X})\otimes_{\unsure{\mathbf{Z}|\mathbf{Q}}}\mathbf{Q}&\text{si }i=2j\deleted{,2j+1}\\
\deleted{\unsure{A^j(\overline{X})}}\ \mathrm{id}&\text{si }i=2j+1,
\end{cases}\\
\Vert&&\text{mais pas \uuline{canoniquement} en }\unsure{k}\ !\\
H^i(\overline{X},\mathbf{Q}(j))&&\inserted{\unsure{fonctoriellement}}
\end{array}
\]
\margin{\uline{N.B.} \unsure{Remplacer} $k$ par $k_n$ revient à multiplier l'isom obtenu par \unsure{$n$} ---}
